% !TEX program = xelatex
% ======================================================================
% 通用中文数学建模论文模板（EasyMathModel 内置默认模板）
% ----------------------------------------------------------------------
% 适用场景：赛制没有提供官方论文模板时，作为 paper/main.tex 的默认骨架。
% 使用方法：
%   1. 把 templates/paper/ 复制为工作区 paper/；
%   2. 把各问正文写入 paper/sections/q1.tex、q2.tex ...；
%   3. 在 refs.bib 中填写真实参考文献（禁止编造）；
%   4. 填写中英文摘要、关键词、标题；
%   5. 编译：xelatex main && bibtex main && xelatex main && xelatex main。
% 重要：如果赛制提供了官方模板（.tex/.docx 模板或排版要求），
% 必须优先使用官方模板；本文件只是没有官方模板时的默认骨架。
% ======================================================================
\documentclass[12pt,a4paper,fontset=fandol]{ctexart}

\usepackage[a4paper,top=2.54cm,bottom=2.54cm,left=3.17cm,right=3.17cm]{geometry}
\usepackage{amsmath,amssymb,amsfonts}
\usepackage{graphicx}
\graphicspath{{figures/}}
\usepackage{booktabs,multirow,longtable,array}
\usepackage{float}
\usepackage{caption,subcaption}
\usepackage{listings}
\usepackage{xcolor}
\usepackage{enumerate}
\usepackage[colorlinks=true,linkcolor=blue,citecolor=blue,urlcolor=blue]{hyperref}

\lstset{
  basicstyle=\ttfamily\footnotesize,
  breaklines=true,
  frame=single,
  numbers=left,
  numberstyle=\tiny\color{gray},
  showstringspaces=false
}

\setlength{\parindent}{2em}
\linespread{1.2}

\begin{document}

% ---------- 标题（按赛制要求匿名时，不写作者信息） ----------
\begin{center}
  {\LARGE\bfseries 论文标题（按赛题重写）}
\end{center}

% ---------- 中文摘要 ----------
\begin{abstract}
  （中文摘要正文：每个子问题独立成段，先写结论后写方法；
  摘要中每个数字必须在正文中原样出现。）
  \vspace{2mm}

  \noindent\textbf{关键词：}关键词一；关键词二；关键词三
\end{abstract}

% ---------- 英文摘要 ----------
\begin{center}
  {\Large\bfseries Abstract}
\end{center}
English abstract goes here. Summarize the main method and the headline
results in one paragraph per subquestion.

\newpage
\tableofcontents
\newpage

% ---------- 正文骨架（由 paper-section-writer 填充） ----------
\section{问题重述}
\input{sections/restatement}

\section{问题分析}
\input{sections/analysis}

\section{模型假设}
\input{sections/assumptions}

\section{符号说明}
% 符号表使用 longtable，保证跨页可断；所有符号与 planning/symbol_table.md 一致。
\input{sections/symbols}

\section{模型的建立与求解}
\subsection{问题一}
\input{sections/q1}
\subsection{问题二}
\input{sections/q2}
\subsection{问题三}
\input{sections/q3}
\subsection{问题四}
\input{sections/q4}
\subsection{问题五}
\input{sections/q5}

\section{灵敏度与稳健性分析}
\input{sections/robustness}

\section{模型评价与推广}
\input{sections/evaluation}

% ---------- 参考文献（BibTeX，真实检索获取） ----------
\bibliographystyle{unsrt}
\bibliography{refs}

% ---------- 附录：代码不进正文 ----------
\appendix
\section{附录：主要程序代码}
% 代码放在附录，用 listings 排版；禁止在正文中出现代码块。
\input{sections/appendix}

\end{document}
